\setcounter{section}{1}\setcounter{theo}{0}
% Official AHL_2025__8__555_0.tex lines 208--214
\begin{theo}\label{threequarters}
The maximal Euler characteristic $\chi(d)$ among all smooth connected cobordisms between the links $T(d,d)$ and $T(2,\lfloor\frac{3}{4}d^2\rfloor)$ satisfies
\[
\lim_{d \to \infty} \frac{\chi(d)}{d^2}=-\frac{1}{4}.
\]
\end{theo}


% Official AHL_2025__8__555_0.tex lines 221--228
\begin{theo}\label{prop:cobdist}
For all non-zero integers $d$ and all integers $N$, the maximal Euler characteristic $\chi(d,N)$ among all smooth connected cobordisms between the links $T(d,d)$ and $T(2,N)$ has value around $-\frac{1}{4}d^2-|\rmN-\signn(d)\frac{3}{4}d^2|$. More precisely,
\[
-\frac{d^2}{4}-\left|\signn(d)\frac{3}{4}d^2-N\right|-4|d|
\leq\chi(d,N)\leq -\frac{d^2}{4}-\left|\signn(d)\frac{3}{4}d^2-N\right|+2|d|.
\]
\end{theo}


% Official AHL_2025__8__555_0.tex lines 254--387
\section{Hirano's curves and the proof of Theorem~\ref{threequarters} }\label{sec:main}

We consider links $L\subset S^3$ smooth non-empty 1-submanifolds and cobordisms between them. For us, a \emph{cobordism} between two links $L_0,L_1\subset S^3$ is a smooth oriented neatly embedded surface $F\subset S^3\times[0,1]$ such that $\partial F=L_0\times\{0\}\cup \bar{L_1}\times\{0\}$, where $\bar{L_1}$ denotes $L_1$ with reversed orientation. While all our cobordisms are smooth, we will add the adjective ``smooth'' in our statements, to avoid any confusion with the topological category that behaves fundamentally differently.

We derive Theorem~\ref{threequarters} (and Theorem~\ref{prop:cobdist}) from the following family of examples.

\begin{lemm}\label{lem:chiexact}
For all integers $m\geq 1$, there exists a connected smooth cobordism of Euler characteristic $-m^2-2m$ between the links $T(2m,2m)\text{ and } T(2,3m^2)$.
\end{lemm}

Lemma~\ref{lem:chiexact} follows from the existence of a particular family of projective algebraic curves: for every integer $m\geq 2$, there exists an irreducible projective algebraic curve in $\C\P^2$ of degree~$2m$ with exactly $N:= 3m$ singularities, all of which are of type~$A_{m-1}$~\cite[Theorem~2]{Hirano_92}; see also~\cite[Theorem~3.2]{GS08}.

\begin{proof}[Proof of Lemma~\ref{lem:chiexact}]
For $m=1$, there even exists a smooth connected cobordism from $T(2,2)$ to $T(2,3)$ with Euler characteristic $-1$. Hence we consider the case of an integer $m\geq 2$. We consider the projective algebraic curve $C\subset \C\P^2$ of degree $d\coloneqq 2m$ given by the irreducible homogeneous polynomial
\[
F= \left(x^{m}+y^{m}+z^{m}\right)^2-4\left(x^{m}y^{m}+y^{m}z^{m}+z^{m}x^{m}\right)\in\C[x,y,z],
\]
which was used by Hirano to prove~\cite[Theorem~2]{Hirano_92}. A calculation reveals that $C$ has $N:= 3m$ singularities, all of which are of type~$A_{m-1}$.

We obtain an affine algebraic curve in $\C^2$ with $N$-many $A_{m-1}$ singularities by removing a generic projective line from $\C\P^2$. We explain this in more detail in the rest of this paragraph. Pick a projective line $L$ (a subvariety $L\subset \C\P$ defined by $ax+by+cz=0$ for some $[a:b:c]\in\C\P^2$) that intersects $C$ transversally (i.e.~if $p\in C\cap L$, then $p$ is a non-singular point of $C$ and the tangent spaces of $C$ and $L$ at $p$ span the entire tangent space of $\C\P^2$ at $p$). Note that by Bézout's theorem~$C\cap L$ consists of $d$ points. We pick a linear transformation $A\in \Gl(3,\C)$ that maps $L$ to the line at infinity, which is defined by $z=0$ and denoted by $\C\P^1$. Now consider $\widetilde{F}=F\circ A^{-1}$, which is a homogeneous polynomial defining the curve $\widetilde{C}$ obtained from applying $A$ to $C$, define $\widetilde{f}\coloneqq \widetilde{F}(x,y,1)$, and take the desired affine algebraic curve in $\C^2=\C\P^2\setminus \C\P^1$ to be $\widetilde{f}^{-1}(0)$.

The fact that $\widetilde{C}$ intersects $\C\P^1$ transversally, implies that the link at infinity of $\widetilde{f}^{-1}(0)$ is the torus link $T(d,d)$, i.e.~$S^3_R\cap \widetilde{f}^{-1}(0)$ has link type $T(d,d)$ for $R>0$ large~enough. Here $S^3_R\subset \C^2$ denotes the $3$-sphere of radius $R$, i.e.~the boundary of the $4$-ball $B^4_R$ of radius $R$ with center the origin. In order to see this, consider a closed regular neighborhood $\nu(\C\P^1)$ of $\C\P^1\subset\C\P^2$. To be concrete, take $\nu(\C\P^1)$ to be the complement of $\intt(B^4_R)\subset\C^2\subset\C\P^2$ for some large $R>0$. Such a neighborhood is diffeomorphic (via some $\phi$) to the total space $E$ of the once-twisted $D^2$-bundle over $S^2\cong \C\P^1$ (the $D^2$-bundle $\pi\colon E\to S^2$ with $\partial E=S^3$). Choosing $\nu(\C\P^1)$ smaller if needed (that is increasing $R$), by transversality we can arrange for the diffeomorphism $\phi\colon \nu(\C\P^1)\to E$ to map $\nu(\C\P^1)\cap \widetilde{C}$ to $d$ fibers $\pi^{-1}(p_1),\pi^{-1}(p_2),\dots,\pi^{-1}(p_d)$ for points $p_1,\dots,p_d\in S^2$. Since $\pi|_{\partial E}\colon \partial E\to S^2$ is the Hopf fibration (since, up to isomorphisms of bundles, there is only one $S^1$-bundle over $S^2$ with total space $S^3$), $T\coloneqq\pi|_{\partial E}^{-1}\{p_1,\dots,p_d\}\subset \partial E=S^3$ is a $T(d,d)$ torus link. Recalling 
${\partial(\nu(\C\P^1))=S^3_R}$, we see that $\phi$ yields a diffeomorphism of pairs between $(S^3_R,S^3_R\cap \widetilde{f}^{-1}(0))=(S^3_R,\phi^{-1}(T))$ and $(S^3, T)$. Hence, as desired, the link at infinity is the torus link $T(d,d)$.

We now use $\widetilde{f}^{-1}(0)$ to find the desired cobordism. For this let $s_1,\dots, s_N$ denote the singular points of $\widetilde{f}^{-1}(0)$ and choose $\epsilon>0$ small enough that $S^3_{\epsilon,s_k}$ (the sphere of radius $\epsilon$ around $s_k$) intersects $\widetilde{f}^{-1}(0)$ transversally in $T(2,m)$. Let $W:= B_R^4\setminus(\bigcup_{1\leq k\leq N}\intt(B^4_{\epsilon,s_k}))$, where $B^4_R$ denotes the close ball with boundary $S^3_R$ and $\intt(B^4_{\epsilon,s_k})$ denotes the open ball with boundary $S^3_{\epsilon,s_k}$. Towards finding our cobordism (a surface $F$ in $S^3\times [0,1]$), we modify $W$ to be diffeomorphic to $S^3\times [0,1]$ by tubing the small boundary spheres together and track what happens to $\widetilde{f}^{-1}(0)$. For this, we pick $N-1$ pairwise disjoint embedded closed arcs $a_k\subset W\cap\widetilde{f}^{-1}(0)$ that start and end on $\partial W\setminus S^3_R=\bigcup_{1\leq k\leq N} S^3_{\epsilon,s_k}$ such that $(\partial W\setminus S^3_R)\cup \bigcup_{1\leq k\leq N-1} a_k$ is connected. We take $\nu(a_k)$ to be a small closed tubular neighborhood of $a_k$ such that its boundary intersects $\widetilde{f}^{-1}(0)$ transversally and the pair $(\nu(a_k), \widetilde{f}^{-1}(0)\cap\nu(a_k))$ is diffeomorphic to $(B^3\times[0,1],([-1,1]\times\{0\}\times\{0\})\times[0,1])$, where $B^3$ denotes the closed unit ball in~$\R^3$. We set $X$ to be the closure in $\C^2$ of $W\setminus\bigcup_{1\leq k\leq N} \nu(a_k)$. After smoothing corners, we pick a diffeomorphism from $X$ to $S^3\times [0,1]$ and let $F\subset S^3\times [0,1]$ denote the image of $X\cap\widetilde{f}^{-1}(0)$ under said diffeomorphism.

Next, we determine the Euler characteristic of $F$ and along the way observe that it is connected. We discuss the case that $m$ is odd. A similar calculation works when $m$ is even and yields the same result. We first note that $C$ (as a topological surface) is connected, as is the case for all irreducible projective curves in $\C\P^2$, and has genus $(d-1)(d-2)/2-N\frac{m-1}{2}$. The latter can for example be seen by noting that a small generic deformation of $C$ is a smooth algebraic curve of degree $d$, which has genus $(d-1)(d-2)/2$, where each of the $N$-many $A_{m-1}$-singularities contributes $\frac{m-1}{2}$ to the genus. $\widetilde{f}^{-1}(0)$ and $\widetilde{f}^{-1}(0)\cap W$ have the same genus as $C$ (since they are obtained by removing discs), and $\widetilde{f}^{-1}(0)\cap X\cong F$ also has the same genus as $C$ and is also connected since it is obtained from $\widetilde{f}^{-1}(0)\cap W$ by removing neighborhoods of embedded arcs that connect different boundary components. Thus, we find
\[
\chi(F)=2-2\frac{(d-1)(d-2)-N(m-1)}{2}-(d+1) =-(d-1)(d-1)+N(m-1),
\]
because $F$ is connected and has $d+1$ boundary components.

By construction, $F$ is a connected smooth cobordism between $T(d,d)$ and a knot $K$ given as the connected sum of $N$-many $T(2,m)$ torus links. Take $F'$ to be a cobordism given by $N-1$ one-handles between $K$ and $T(2,Nm)$; in particular, $F'$~is connected and $\chi(F')=-N+1$.

Composing the two cobordisms $F$ and $F'$, we find a connected smooth cobordism from $T(d,d)$ to $T(2,Nm)$ with Euler characteristic
\[
-(d-1)(d-1)+N(m-1)-N+1=-(d-1)(d-1)+(m-2)N+1=-m^2-2m.\qedhere
\]
\end{proof}

\begin{proof}[Proof of Theorem~\ref{threequarters}]
We fix an integer $d\geq 2$ and consider the largest integer~$m$ such that $d\geq D\coloneqq 2m$. In other words, $D=d$ if $d$ is even and $D=d-1$ if $d$ is odd. We construct the desired cobordism as the composition of three cobordisms, $F$, $G$, and $F'$. The key contribution is $H$, which is given by Lemma~\ref{lem:chiexact}, while $F$ and~$F'$ are of a more cosmetic nature that help adapt the interesting cobordism $H$ ($H$ stands for Hirano) to start and end at the right links. We describe all three in detail.

Let $F$ be a smooth cobordism from $T(d,d)$ to $T(D,D)$ with
\[
\chi(F)=(D-1)^2-(d-1)^2=
\begin{cases}
0&\text{if $d$ is even}\\
-4m+1&\text{if $d$ is odd}
\end{cases}
\]
as follows: $F$ is the identity cobordism if $d$ is even, i.e.~$F=T(d,d)\times[0,1]$, while, for $d$ odd, $F$ is any connected cobordism. We expand on the existence of such an~´$F$ for $d$ odd in the rest of this paragraph. For example, $F$ can be obtained by a complex algebraic curve $C\in \C^2$ intersected with $B^4\setminus \intt(B_{\frac{1}{2}})$, where $C$ is chosen to intersect $S^3$ and $S^3_{\frac{1}{2}}$ transversally in $T(d,d)$ and $T(D,D)$, respectively. (See for example Remark~\ref{rem:T(d,d)atinfty} for how such a curve $C$ can be produced.) The local Thom conjecture implies that the Euler characteristic is as claimed. Here is a different, explicit construction that builds a handle decomposition of a cobordism as desired. Consider the following sequence of $4m=2D=2(d-1)$ positive $d$-braids:
\begin{align*}
\sigma_1\sigma_2,\dots,\sigma_{d-2}\sigma_{d-1}(\sigma_1\sigma_2,\dots,\sigma_{d-2})^{d-2}, (\sigma_1\sigma_2,\dots,\sigma_{d-2}\sigma_{d-1})^2(\sigma_1\sigma_2,\dots,\sigma_{d-2})^{d-3},\\
(\sigma_1\sigma_2,\dots,\sigma_{d-2}\sigma_{d-1})^3(\sigma_1\sigma_2,\dots,\sigma_{d-2})^{d-4},
\dots,
(\sigma_1\sigma_2,\dots,\sigma_{d-2}\sigma_{d-1})^{d-1},\\ (\sigma_1\sigma_2,\dots,\sigma_{d-2}\sigma_{d-1})^{d-1}\sigma_1,
(\sigma_1\sigma_2,\dots,\sigma_{d-2}\sigma_{d-1})^{d-1}\sigma_1\sigma_2,
\dots,\\
(\sigma_1\sigma_2,\dots,
\sigma_{d-2}\sigma_{d-1})^{d-1}\sigma_1\sigma_2,\dots,
\sigma_{d-2},\\
(\sigma_1\sigma_2,\dots,\sigma_{d-2}\sigma_{d-1})^{d-1}\sigma_1\sigma_2,\dots,
\sigma_{d-1}\\
=(\sigma_1\sigma_2,\dots,\sigma_{d-2}\sigma_{d-1})^{d}.
\end{align*}
The first braid has as its braid closure $T(D,D)$ and the last braid has $T(d,d)$ as its braid closure. Each differs from its predecessor by one more generator $\sigma_i$, hence their closures are related by a cobordism consisting of one $1$-handle (in particular a cobordism of Euler characteristic $-1$). Composing these $4m-1$ cobordism yields the desired cobordism $F$. To see that $F$ is connected, note that the $(d-1)^{\rm th}$ braid has closure $T(d, d-1)$ (hence is connected) and by construction (only $1$-handles) every point in $F$ has at most as many connected component as this level set.

Let $F'$ be a smooth cobordism from $T(2,3m^2)$ to $T(2,\lfloor \frac{3}{4}d^2\rfloor)$ with
\[
\chi(F')=-\left\lfloor 3d^2/4\right\rfloor+3m^2=
\begin{cases}
0&\text{ if $d$ is even}\\
-3m&\text{ if $d$ is odd}
\end{cases}
\]
as follows: if $d$ is even, $F'$ is the identity cobordism, while if $d$ is odd, $F'$ is any connected cobordism. Similar to $F$, if $d$ is odd, such an $F'$ can for example be given by an algebraic curve or an explicit construction as follows. Consider the following sequence of $3m+1=\lfloor 3d^2/4\rfloor-3m^2$ positive $2$-braids:
\[
\sigma_1^{3m^2},
\sigma_1^{3m^2+1},
\dots,
\sigma_1^{\left\lfloor 3d^2/4\right\rfloor},
\]
for which the closure of the first and last braid are $T(2,3m^2)$ and $T(2,\lfloor 3d^2/4\rfloor)$, respectively. The closure of any two consecutive braids are related by a cobordism given by a one handle that is connected. Since one of the closures has two components while the other has one, all these $3m$ cobordisms are diffeomorphic to a $3$-holed sphere; in particular they are connected and have Euler characteristic $-1$. Hence, their composition yields a connected cobordism $F'$ of Euler characteristic $-3m$ as desired.

Composing the following three cobordisms
\begin{itemize}
\item $F$ from $T(d,d)$ to $T(D,D)$,
\item a cobordism $H$ from $T(D,D)$ to $T(2,3m^2)$ as guaranteed to exist by Lem\-ma~\ref{lem:chiexact},
\item and $F'$ from $T(2,3m^2)$ to $T(2,\lfloor \frac{3}{4}d^2\rfloor)$,
\end{itemize}
we find a connected cobordism $G$ between $T(d,d)$ and $T(2,\lfloor\frac{3}{4}d^2\rfloor)$ with
\begin{align*}
\chi(G)&= \chi(F)+\chi(H)+\chi(F')
\\&=
\begin{cases}
0-m^2-2m+0=-m^2-2m&\text{ if $d$ is even}\\
-4m+1-m^2-2m-3m=-m^2-9m+1&\text{ if $d$ is odd}
\end{cases}
\end{align*}
Therefore, we have
\[
\chi(d)\geq
\begin{cases}
-m^2-2m=-\frac{d^2}{4}-d&\text{ if $d$ is even}\\-m^2-9m+1=-\frac{d^2}{4}-4d+5+\frac{1}{4}&\text{ if $d$ is odd}
\end{cases}
\geq-\frac{d^2}{4}-4d.
\]

To find an upper bound on $\chi(d)$, we employ Murasugi's signature obstruction on the Euler characteristic of cobordism between links~\cite{Murasugi_OnACertainNumericalInvariant}. Using
\begin{equation}\label{eq:sigfortorusknots}
\sigma(T(d,d))=-\left\lfloor\frac{d^2-1}{2}\right\rfloor\text{ and }\sigma\left(T\left(2,k\right)\right)=-k+1\quad\text{\cite[Theorem~5.2]{GLM}},
\end{equation}
for all positive integers $d$ and $k$, we find
\[
\chi(d)\overset{\text{\cite{Murasugi_OnACertainNumericalInvariant}}}{\leq}
\sigma\left(T\left(2,\left\lfloor\frac{3}{4}d^2\right\rfloor\right)\right)-\sigma(T(d,d))\overset{\text{\eqref{eq:sigfortorusknots}}}{=}\left\lfloor\frac{d^2-1}{2}\right\rfloor-\left\lfloor\frac{3}{4}d^2\right\rfloor+1\leq -\left\lfloor\frac{d^2}{4}\right\rfloor+1.
\]

In conclusion, we have shown
\begin{equation}\label{eq:boundsonchi(d)}
-\frac{d^2}{4}-4d\leq\chi(d)\leq-\left\lfloor\frac{d^2}{4}\right\rfloor+1;
\end{equation}
in particular, this establishes $\lim_{d \to \infty} \frac{\chi(d)}{d^2}=-\frac{1}{4}$.
\end{proof}

We end this section with a remark about rearranging curves in $\C^2$ such that their link at infinity is $T(d,d)$ without changing the singularity at the origin, which we have used in the first paragraph of the introduction. The argument is very similar to the one from the second paragraph of the proof of Lemma~\ref{lem:chiexact}.
\begin{rema}\label{rem:T(d,d)atinfty}
Let $C\subset\C^2$ be a reduced\footnote{Reduced simply amounts to the defining polynomial $f$ being square-free.} algebraic curve of degree $d$ and fix some $p\in C$. We claim by adding a small degree $d$ polynomial to a defining square-free polynomial $f$ of $C$, we can change $C$ to a reduced algebraic curve $\widetilde{C}$ with the same singularity at $p$ such that the link at infinity of $\widetilde{C}$ is $T(d,d)$. This can be done similarly as argued in the second paragraph of the proof of Lemma~\ref{lem:chiexact}. Compare also with~\cite[Section~2]{Feller_15_MinCobBetweenTorusknots}.

Considering the closure in $\C\P^2$ (by homogenizing to a 3-variable polynomial $F$ of degree $d$ with $F(x,y,1)=f(x,y)$), choosing a generic projective line, and then composing the defining equation $F$ with a linear transformation of $\C\P^2$ that maps this line to the line at infinity $\{[x:y:z]\mid z=1\}$ and fixes $p$, we find a new defining equation $\widetilde{F}$ such that setting $\widetilde{f}(x,y)=\widetilde{F}(x,y,1)$ yields the defining equation of an algebraic curve $\widetilde{C}$ as desired. To guarantee that the coefficients between $f$ and $\widetilde{f}$ vary little, choose the generic projective line to be given by an equation $ax+by+cz=0$ with $a$ and $b$ close to $0$ and $c$ close to $1$ and choose the linear transformation close to the identity (say as an element in $\Gl_3(\C)$). Finally, we note that a further small change of the coefficients of $\widetilde{f}$ (respectively $\widetilde{F}$) does not change the fact that the link at infinity is $T(d,d)$ since the corresponding projective curve transversally intersects the line at infinity.

In the second paragraph of the introduction, we further wanted a smooth algebraic curve. For this we note that adding a generic constant arranges that $\widetilde{f}^{-1}(0)$ is smooth. Choosing the constant small (compared to an $\epsilon$ with $S^3_{\epsilon'}\pitchfork \widetilde{f}^{-1}(0)$ being a torus link of type $T(2,n+1)$ for all $0<\epsilon'\leq\epsilon$) assures that $\widetilde{f}^{-1}(0) \cap S^3_\epsilon$ remains a torus link of type $T(2,n+1)$. Also, as argued at the end of the last paragraph, the link at infinity remains $T(d,d)$ if the constant is chosen sufficiently small.
\end{rema}

% Official AHL_2025__8__555_0.tex lines 389--434
\section{Cobordism distance between torus links of type \texorpdfstring{$T(d,d)$}{T(d,d)} and \texorpdfstring{$T(2,N)$}{T(2,N)}}\label{sec:asymdist}

Theorem~\ref{threequarters} and its proof combined with the link signature bound on the Euler characteristic of cobordisms and the resolution of the local Thom conjecture allows to determine the smooth cobordism distance between the link $T(d,d)$ and all torus links $T(2,N)$ up to an error of at most $3d$. This is the content of Theorem~\ref{prop:cobdist}, which we prove below, after a comment on the cobordism distance.

One may define the cobordism distance between two links $L_0$ and $L_1$ as minus the maximal Euler characteristic of all smooth cobordisms $F$ between $L_0$ and $L_1$ for which each component of $F$ has boundary components in both $S^3\times\{0\}$ and $S^3\times\{1\}$. With this definition the cobordism distance is an integer valued metric on the set of smooth concordance classes of links, justifying the name ``distance''. Alternatively, the first author considered a variation when investigating cobordism distance between positive torus links using scissor equivalences; see~\cite{Baader_ScissorEq}. In case of two non-isotopic non-trivial positive torus links, whatever variation of the definition is chosen, the cobordism distance is realized by a connected smooth cobordism. For clarity of presentation, we abstain from using the term cobordism distance in our statements and instead consider connected smooth cobordisms, as done in the statement of Theorem's~\ref{threequarters} and~\ref{prop:cobdist}.

\begin{proof}[Proof of Theorem~\ref{prop:cobdist}]
Without loss of generality, take $d$ to be positive. In fact, we consider the case when $d\geq 2$, since the case $d=1$ is immediate from the local Thom conjecture; compare~\eqref{eq:locthom} below.

We first discuss the case $N\geq0$. Using a cobordism between $T(d,d)$ and $T(2,\lfloor \frac{3}{4}d^2\rfloor)$ that realizes $\chi(d)$ and composing it with a connected cobordism between $T(2,\lfloor \frac{3}{4}d^2\rfloor)$ to $T(2,N)$ of Euler characteristic $-|\lfloor \frac{3}{4}d^2\rfloor-N|$, yields $-|\lfloor \frac{3}{4}d^2\rfloor-N|+\chi(d)\leq\chi(d,N)$. Combined with
\[
-\left|\left\lfloor \frac{3}{4}d^2\right\rfloor-N\right|-\frac{d^2}{4}-4d
\overset{\text{\eqref{eq:boundsonchi(d)}}}{\leq}-\left|\left\lfloor \frac{3}{4}d^2\right\rfloor-N\right|+\chi(d),
\]
we find the desired lower bound
\[
-\frac{d^2}{4}-\left|\frac{3}{4}d^2-N\right|-4d
\leq\chi(d,N).
\]

For the upper bound, in case $N\geq \frac{3}{4}d^2$, we recall the signature bound
\[
\chi(d,N)\overset{\text{\cite{Murasugi_OnACertainNumericalInvariant}}}{\leq} -\sigma(T(d,d))+\sigma(T(2,N))\overset{\eqref{eq:sigfortorusknots}}{=} \left\lfloor\frac{d^2-1}{2}\right\rfloor-N+1{\leq} -\left|-\frac{d^2}{2}+N\right|+\frac{1}{2}
\]
and apply $|-\frac{d^2}{2}+N|=|\frac{d^2}{4}+(-\frac{3}{4}d^2+N)|= \frac{d^2}{4}+|-\frac{3}{4}d^2+N|$ to find
\[
\chi(d,N)\leq -\frac{d^2}{4}-\left|-\frac{3}{4}d^2+N\right|+\frac{1}{2},
\]
as desired. If instead, $0\leq N< \frac{3}{4}d^2$, we use the following triangle inequality for the cobordism distance $\chi(d,N)\leq \chi(d,1) - \chi(1,N)$ in combination with the following consequence of the local Thom conjecture~\cite[Corollary~1.3]{KronheimerMrowka_Gaugetheoryforemb}:
\begin{equation}\label{eq:locthom}
\chi(d,1)=-(|d|-1)^2\text{ and } \chi(1,N)=-\left||N|-1\right|,
\end{equation}
for all integers $N$ and non-zero integers $d$. We find
\begin{align*}
\chi(d,N)&\leq-(d-1)^2+(N-1)
\\&= -d^2+N+2d-2
\\&=-\frac{d^2}{4}-\left|\frac{3}{4}d^2-N\right|+2d-2,
\end{align*}
where we combined the triangle inequality and~\eqref{eq:locthom} to see the inequality.

Finally, if $N\leq -1$, then
\[
\chi(d,N)= \chi(d,1)+\chi(1,N)\overset{\text{\eqref{eq:locthom}}}{=}-(d-1)^2+N+1 =-\frac{d^2}{4}-\left|\frac{3}{4}d^2-N\right|+2d,
\]
where the first equality is a consequence of the local Thom conjecture~\cite[Corollary~1.3]{KronheimerMrowka_Gaugetheoryforemb}.
\end{proof}

